\section{已有数值证据与最小验证设计}
\subsection{精确反馈的保存结果}
本节前半部分只读取旧稿已保存数据，未重跑实验。
原二维模型为
\[
A_{\rm S}=\begin{bmatrix}.82&.18\\-.10&.72\end{bmatrix},\qquad
A_{\rm U}=\begin{bmatrix}1.04&.12\\.03&.91\end{bmatrix},
\]
\[
P_0=\begin{bmatrix}.8&.12\\.12&.6\end{bmatrix},\quad
Q=\begin{bmatrix}.05&.01\\.01&.04\end{bmatrix},\quad
R=\begin{bmatrix}.2&.03\\.03&.3\end{bmatrix},\quad
Z=\begin{bmatrix}2&.25\\.25&1.5\end{bmatrix}.
\]
目标$k=4$、$b=0.1$。下面两类合法策略的参数沿用旧稿：
持续定向策略，以及目标拍随机符号、下一拍补偿的原两拍策略。
本稿第3节已采用更简单的单拍下界证明；旧两拍结果仍可作为合法见证，
但其数值常数不移植给新的单拍证法。

\begin{table}[ht]
\centering\small
\caption{旧数据中的真实目标损伤分解（精确反馈）}
\begin{tabular}{llrrr}
\toprule
模型 & 原保存策略 & 偏差平方 & 协方差变化 & 真实损伤$J$\\
\midrule
S & 持续定向 & .001315759 & -.000020445 & .001295314\\
U & 持续定向 & .003252364 & -.000120021 & .003132343\\
S & 原两拍符号 & .000321850 & 约0 & .000321850\\
U & 原两拍符号 & .000573797 & 约0 & .000573797\\
\bottomrule
\end{tabular}
\end{table}
持续策略的协方差变化为负，说明只报告条件偏差平方会高估这些例子的真实损伤。
原两拍策略在目标之前无动作，目标协方差分布不变，所以该项为零。
原档案使用事件树计算并以每种组合40万次独立重置结果交叉检查，
本整合不产生额外重复证据。

\begin{table}[ht]
\centering\small
\caption{不同种类的损伤界不能混用（旧两拍策略，$b=0.1$）}
\begin{tabular}{lrr}
\toprule
指标 & S & U\\
\midrule
真实损伤 & .000321850 & .000573797\\
总体统计证书 & .000457526 & .000902386\\
总体证书/真实损伤 & 1.42155 & 1.57266\\
双预算上界/原构造下界，TV线性区 & 783.639 & 359.136\\
$N=10^4$、已知参考的总体中心置信界示意 & .084711 & .072733\\
\bottomrule
\end{tabular}
\end{table}
最后一行使用旧Hoeffding公式并把经验均值置于总体均值，
不是新采样所得认证结果。总体证书在这两个例子中较接近真实损伤，
有限样本余量却仍很大。双预算曲线只是全策略上界与一条合法构造下界，
不代表真实最优值已经被数值求解。原事件串TV也不替代完整连续日志TV；
后者预算可行性由解析联合核界提供。

\subsection{新增机制验证：设计已固定，尚未执行}
\noindent\fbox{\parbox{.94\linewidth}{\small 本小节为验证设计，不包含量化性能结果。
不得把下述参数表、理论趋势或预期一致性写成已完成实验。}}

量化定理目前限标量，故新增验证使用两个独立标量系统S1/U1，
$A=0.82/1.04$，$m_0=0$、$P_0=0.8$、$Q=0.05$、$R=0.2$、$Z=2$，目标$k=4$。
这些数值来源于旧参数的首坐标系数，但删除耦合后不是旧二维系统的边缘。
固定$b=b_{\rm ref}=0.1$。在已证明统一概率范围内确定三个精度：
\begin{center}\small
\begin{tabular}{lrrr}
\toprule
模型 & $\delta_1$ & $\delta_2$ & $\delta_3$\\
\midrule
S1 & .00290706 & .00581412 & .01162823\\
U1 & .00556494 & .01112988 & .02225976\\
\bottomrule
\end{tabular}
\end{center}
这些值仅由有限协方差标签的确定性常数计算得到，规则为
\[
\delta_{\rm cap}=\min\{0.1,\delta_*/2,\alpha_{\max}b/2\},\qquad
\delta_j=\delta_{\rm cap}\{1/4,1/2,1\}.
\]
精度列表没有按性能筛选。

每模型预定2048条独立基础物理与触发样本，A/B及三个精度共享样本。
每个样本平均两个私有符号分支，统计独立单位仍是基础样本。
A采用精确中心；B采用量化中心和一致静默更新。
二者都在本运行目标拍已真实发送后，依自己的合法前史按中央格条件选择同形式的有界见证。
C是B的额外观察视图，不能重跑成另一条状态轨迹。

必须保存真实目标MSE、同协议健康基线、通信次数、幅值/实际能量、
事件和量化码、P标签、接收载荷和未量化预测，以及相应的来源权限。
重点检验B的事件/反馈路径一致性、正损伤恒等式及C的解析矩差。
不要求A与B实际通信成本恰相等，也不据此排名策略性能。
TV同律由数学证明给出；样本中的零不一致仅检查代码实现。
探索性直方图TV不作为完整日志总体TV的估计。

待完成上述冻结验证后，本节拟用一张精度--损伤/正常代价图和一张信息对照表补全。
若有限样本无法区分A的很小事件概率差，保留区间并用解析概率解释，不增加种子寻找显著性。
当前阶段不运行一般维数量化、固定比特或三预算新实验。
